%% Écrit par outils/glossaire.py — ne pas modifier à la main
\annexe[glossaire]{Glossary}{The programmer's vocabulary, from A to Z}
\begin{glossaire}
\glolettre{A}
\gloterme{augmented assignment}{An instruction such as \py{x += 3} that updates a variable using its current value. It is a shorter way of writing \py{x = x + 3}.}{3}
\glolettre{C}
\gloterme{comment}{A note in the code, written after a \py{\#} sign. Python ignores it: it is there to help humans understand the program.}{1}
\glolettre{E}
\gloterme{expression}{A piece of code that Python evaluates to produce a value: \py{7 + 5}, \py{3 * 4 - 1}, or simply \py{42}. An \textbf{operator} is a sign that combines values, such as \py{+} or \py{*}.}{2}
\glolettre{F}
\gloterme{function}{A named tool that performs a task. You call a function by writing its name followed by parentheses: \py{print("Hi")}. The values inside the parentheses are its \textbf{arguments}.}{1}
\glolettre{I}
\gloterme{IDLE}{The editor installed with Python. It has two kinds of windows: the \textbf{Shell}, to talk to Python one instruction at a time, and the \textbf{editor}, to write and save complete programs.}{1}
\gloterme{interpreter}{The program that reads Python code and carries it out, instruction by instruction. Installing \enquote{Python} on a computer means installing the Python interpreter.}{1}
\glolettre{P}
\gloterme{program}{A list of instructions that a computer carries out, one after the other, to perform a task. Each instruction is a simple, precise order: display a message, add two numbers, save a file\ldots}{1}
\gloterme{programming language}{A language with strict rules, designed to write instructions that a computer can carry out. The rules that say how instructions must be written are called the \textbf{syntax} of the language.}{1}
\glolettre{R}
\gloterme{remainder}{What is left over after a floor division. In Python, \py{a \% b} is the remainder of the division of \py{a} by \py{b}, and \py{a // b} its quotient. The sign \py{\%} is also called \enquote{modulo}.}{2}
\glolettre{T}
\gloterme{type}{The kind of a value, which decides what you can do with it. \py{int} is the type of whole numbers, \py{float} the type of decimal numbers. The function \py{type()} gives the type of a value.}{2}
\glolettre{V}
\gloterme{variable}{A name that refers to a value stored in the computer's memory. You create a variable by \textbf{assigning} it a value with the sign \py{=}: \py{price = 4.5}.}{3}
\end{glossaire}
